\input{preambulo.tex}

%----------------------------------------------------------------------------------------
%	TÍTULO Y DATOS DEL ALUMNO
%----------------------------------------------------------------------------------------

\title{
\vspace{10ex}
\normalfont \normalsize
\huge \textbf{Prediseño de un Sistema de Big Data}
}
\author{Francisco Javier Sueza Rodríguez}
\date{\normalsize\today}

\glsenablehyper
\makeglossaries
\loadglsentries{glos}

%----------------------------------------------------------------------------------------
%                                     DOCUMENTO
%----------------------------------------------------------------------------------------
\begin{document}


\maketitle

\thispagestyle{empty}

\vspace{75ex}

\begin{center}
    \begin{tabular}{l l}
        \textbf{Centro}: & IES Zaidin-Vergeles \\
        \textbf{Curso}: & Inteligencia Artificial y Big Data\\
        \textbf{Módulo}: & Sistemas de Big Data \\
    \end{tabular}
\end{center}

\newpage

\begingroup
\hypersetup{linkcolor=black}
\tableofcontents
\endgroup

\newpage

\begingroup
\hypersetup{linkcolor=black}
\listoffigures             
\endgroup

\newpage

\section{Caso Práctico}
\textbf{Logística Nova XXI} es una empresa de transporte y distribución que quiere explorar soluciones Big Data para modernizar su gestión. Sus características y necesidades son las siguientes:

\begin{itemize}
	\item Hay distintas fuentes:
	\begin{itemize}
		\item Externas a la empresa: datos abiertos de tráfico, climatología y precios de carburantes, actualizados varias veces por semana
		\item Internas a la empresa: telemetría en tiempo real de la flota de camiones (GPS, consumo, sensores de mantenimiento) y transacciones de envíos de clientes.
	\end{itemize}
	
	\item La cantidad de datos actualmente es de aproximadamente 50 TB de datos históricos, con previsión de +15 TB/mes por incorporación de nuevos sensores.
	\item Requisitos de negocio:
	\begin{itemize}
		\item Los clientes deben poder consultar en línea el estado de sus envíos (transacciones diarias).
		\item El departamento de operaciones quiere alertas en tiempo real para incidencias (retrasos, averías, desvíos de ruta).
		\item La dirección general necesita un informe semanal con métricas descriptivas (costes de combustible, puntualidad, utilización de la flota) con datos cerrados el domingo.
		\item El área de planificación quiere un modelo predictivo que anticipe retrasos probables en rutas según datos históricos de tráfico y clima.
	\end{itemize}
\end{itemize}

\section{Proyecto de Big Data}
En este apartado vamos a explicar el proyecto de Big Data que se ha diseñado para solucionar el problema que se ha planteado. Se responderán a todas las cuestiones que se plantean en el enunciado y se expondrá el diseño realizado.


\subsection{Identificación de las 5V} 
En este primer ejercicio vamos a identificar las 5Vs en el caso propuesto, justificando la necesidad de emplear un sistema de Big Data como solución.

\begin{itemize}
	\item \textbf{Volumen}
	
	 En nuestro caso el volumen de datos que se especifica en el caso práctico es muy grande, no solo hablamos de más de \textbf{15 TB de datos mensuales}, teniendo en cuenta además que en la mayoría de los casos los registros van ser registros que se pueden representar en texto, es muchísima información. Pero es que además debemos tener en cuenta los \textbf{50 TB acumulados de datos históricos}, que es una cantidad enorme de datos. Si se mantiene la previsión de generación de datos, en un año se habrán incorporado más de 180 TB de información.
	 
	 Esto justifica la implementación de un sistema Big Data, ya que el volumen que ya existe de datos y el que se generará es muy grande.
	
	\item \textbf{Velocidad}
	
	La velocidad es crucial en algunos de los aspecto de nuestro caso, la \textbf{telemetría de los camiones}, los \textbf{datos de tráfico} o la \textbf{climatología} son datos que deben procesarse rápidamente, algunos en tiempo real, para tomar decisiones sobre las rutas que se deben tomar o comprobar el estado de los camiones. Además, los clientes deben consultar el \textbf{estado de la entrega} en tiempo real, por lo que los datos de \textbf{posicionamiento} deben igualmente gestionarse en \textbf{tiempo real}.
	
	Esta velocidad de procesamiento, y de esa cantidad de datos, solo se puede lograr implementado un sistema de procesamiento distribuido.
	
	\item \textbf{Variedad}
	
	En este caso, tenemos una \textbf{gran variedad de fuentes} y de \textbf{formatos}. Tenemos datos \textbf{estructurados} como pueden ser la información de los clientes o el registro de operaciones o datos en formato \textbf{JSON} procedente de sensores, dispositivos de GPS, así como datos de externos que pueden venir en diferentes formatos, como la \textbf{información de tráfico} o \textbf{meteorológica}, \textbf{precios de carburantes},... entre otros. 
	
	En resumen, tenemos una gran variedad de fuentes de datos que pueden venir en diferentes formatos, y que necesitan ser adaptados e integrados conjuntamente para representar la situación con coherencia.
	
	\item \textbf{Veracidad}
	
	En el caso que estamos tratando la veracidad la información es importante. Hay que tener en cuenta que muchos de los datos que podemos recolectar pueden tener errores, ya sea por \textbf{sensores defectuosos}, \textbf{mediciones incorrectas} del consumo de combustible o de estado mecánico de la flota. Además, los datos externos,como los del tráfico o la meteorología pueden tener \textbf{diferentes niveles de precisión}. 
	
	Teniendo en cuenta que estos datos se van emplear para alimentar el modelo predictivo, es importante implementar mecanismos de calidad que se encarguen de verificar la corrección de los datos y de limpiarlos o adaptarlos, en caso de que sea necesario.
	
	\item \textbf{Valor}
	
	Por último hablamos del valor, y en nuestro caso, el valor que nos proporcionan estos datos es una \textbf{mejora generalizada} del funcionamiento de la empresa. Nos permiten tomar decisiones logísticas más eficientes, mejorando la planificación de rutas. Nos permitirán también poder ofrecer servicios al cliente de gran valor, como el seguimiento en tiempo real de sus paquetes. Además, permitirán la reducción de coste al tener un mejor planificación y ayuda en la toma de decisiones estratégicas de negocio.
\end{itemize}

\subsection{Propuesta de Arquitectura Big Data}
En este apartado se va a realizar el diseño del sistema de Big Data que cubra las necesidades mencionadas en caso práctico, que podemos ver de forma resumida en la captura que se muestra a continuación. 

Como nota, \textbf{no se ha incluido} el modelo de \textbf{Machine Learning} en el diagrama, por falta de espacio, aunque se realizará su explicación en la actividad relacionado. Tampoco se ha incluido la \textbf{capa de gobierno}, ya que esta es una capa transversal que afecta a todas, pero si se va a explicar su funcionamiento en este ejercicio.

A continuación de la captura se va a explicar cada capa y lo elementos que la van a componer.

\begin{figure}[H]
	\centering
	\includegraphics[scale=0.60]{arquitectura-bigdata.png}
	\caption{Arquitectura del sistema de Big Data}
\end{figure}

\begin{itemize}
	\item \textbf{Capa de Ingesta} 
	
	En esta capa debemos de tener un software que sea capaz de gestionar los diferentes eventos y mensajes en los que vamos a obtener la información de las diferentes fuentes. Por eso un \textbf{gestor de eventos o mensajes} como puede ser Kafko o RabbitMQ no ayudará a obtener la información y almacenarla en el sistema de almacenamiento adecuado, que veremos a continuación. Esta capa es \textbf{necesaria} para poder distribuir los datos en los sistemas de almacenamiento indicado.
	
	\item \textbf{Capa de Almacenamiento}        
	
	En esta capa tendremos dos almacenamientos diferentes. Por un lado tendremos un \textbf{data lake}, que se encargará de ir almacenando los \textbf{datos históricos} en un sistema de almacenamiento distribuido, como puede ser HDFS, que nos permitiría almacenar grandes cantidades de datos, además de la posibilidad de escalar de forma horizontal. Además, tendremos una \textbf{base de datos relacional} para almacenar la información de los envíos, permitiendo además la realización de consultar rápidas por parte de los usuarios en la página web. Esta capa es \textbf{necesaria} para poder almacenar la información antes de su procesamiento.
	
	\item \textbf{Capa de Procesamiento}   
	
	En esta capa tenemos que procesar la información de diferentes formas. Por un lado, \textbf{los datos históricos} se \textbf{procesarán por lotes}, ya que son grandes cantidades de datos y cuyo resultado no necesita ser tan inmediato. Además, se necesitará \textbf{procesamiento en streaming} para los datos de los pedidos, así como la telemetría de los vehículos y la consultas de los usuarios. También se necesitará realizar la limpieza de los datos, \textbf{mediante procedimientos ETL}.  Podemos usar varias herramientas aunque algunas, como Spache Sparks nos permitirá realizar todo esto. Esta capa es fundamenta para el análisis de los datos y la extracción de valor.
	
	\item \textbf{Data Warehouse}
	
	Este capa, que podría haberse incluido en la de análisis pero se ha decidido poner por separado, se \textbf{encargará de almacenar y preparar} los \textbf{datos históricos} para su análisis. Si no se ha dado en la capa anterior, en esta también se puede implementar el proceso ETL aunque, como hemos comentado, solo afectará a los datos históricos.
	
	\item \textbf{Capa de Análisis}
	
	En esta capa tenemos que distinguir entre componentes diferentes dependiendo de que parte de los datos se encarguen de procesar y gestionar:
	
	\begin{itemize}
		\item \textbf{Inteligencia de Negocio e Informes}: en este apartado estarán las aplicaciones que se encargarán de gestionar los datos de warehouse y realizar análisis de los datos históricos, generando informes que ayuden a la toma de decisiones o a las predicciones.
		
		\item \textbf{Sistemas de Alertas e Información en Tiempo Real}: en este caso los datos se analizarán y usarán directamente desde la capa de procesamiento, donde se emplearán en un sistemas de alertas, para llevar registro de la situación de los envíos y los vehículos, y una interfaz web (o app movil) que será donde los usuarios podrán interactuar con la información de sus pedidos.
	\end{itemize}
	
	\item \textbf{Capa de Gobierno y Seguridad}
	
	Esta capa, que no aparece en el esquema, afecta a todas las capas, ya que se debe encargar de que los datos tiene la calidad suficiente, de los procesos de autenticación y autorización para evitar accesos no permitidos a los datos, de la auditoría y trazabilidad de los datos, documentación del sistema, ...etc. Hay una gran variedad de herramientas y políticas que se deben implementar para cubrir los diferentes aspectos de los que se encarga esta capa. 
\end{itemize}

\subsection{Justificación de Decisiones Clave}

En esta sección se van a explicar el porque se han tomado ciertas decisiones sobre la arquitectura de la aplicación en lo referente a los siguientes conceptos:

\begin{itemize}
	\item \textbf{Escalabilidad}: en nuestro caso se optará por una \textbf{escalabilidad horizontal}, añadiendo nuevos nodos al clúster cuando sea necesario ya que el volumen de datos a procesar va a ir aumentado de forma muy rápida. Respecto al procesamiento, los datos en tiempo real o streaming requieren de un procesamiento muy rápido, por lo que podemos ir ampliando la red de nodos en caso de que estemos cortos de potencia de cómputo.
	
	\item \textbf{ACID vs BASE}: en nuestro caso las propiedades \textbf{ACID} se aplicarían a las \textbf{operaciones transaccionales}, es importante que datos sobre incidencias de envíos, estado de la entrega, o la información relacionada con la facturación se asegure de que las transacciones se completan correctamente y donde la integridad de los datos es fundamental. Por otro lado, se podrían aplicar las propiedades \textbf{BASE} por un lado a datos de telemetría, donde prima su disponibilidad aunque puedan tener un margen de error, así como a los datos de meteorología. Además, se aplicarían a los datos históricos de la aplicación, que estarían almacenados en el data warehouse. 
	
	\item \textbf{OLAP vs OLTP}: al igual que en el caso anterior, incorporaríamos ambos sistemas, ya que cubren aspectos diferentes. Por un lado el sistema \textbf{OLTP} (OnLine Transaction Processing) es necesario para las operaciones transaccionales, las que se van a producir a diario en la empresa especialmente relacionadas con los envíos, gestión de incidencias, estado de la entrega, ...etc, usando para ellos las bases de datos relacionales de la capa de almacenamiento. Por otro lado, el sistema \textbf{OLAP} (OnLine Analytical Processing) que emplearemos sobre todo para realizar las consultas y el análisis sobre los históricos de los datos, aprovechado las capacidades de data warehouse.
	
	\item \textbf{Data Warehouse}: aunque ya lo hemos comentado en los puntos anteriores, \textbf{si sería conveniente} incorporar un Data Warehouse, especialmente para la parte que corresponde a los datos históricos y su procesamiento, ya que facilitaría las consultas analíticas y sería parte clave del sistema OLAP.
\end{itemize} 


\subsection{Estrategias de Procesamiento}

En nuestro sistemas se van usar principalmente 2 estrategias de procesamiento que responden a necesidades diferentes:

\begin{itemize}
	\item \textbf{Procesamiento por Lotes}: es tipo de procesamiento se empleará, principalmente, con los datos que se almacena en el Data Warehouse, ya que tenemos cantidades grandes de datos donde la respuesta no necesita ser tan "inmediata" como  con otros datos de nuestro sistema. Los datos con este tipo de procesamiento se emplearán para realizar análisis de inteligencia de negocios, así como para el entrenamiento del modelo predictivo, aunque de esto hablaremos más adelante. En definitiva, se usarán con los datos históricos.
	
	\item \textbf{Procesamiento en Streaming}: ese procesamiento se empleará en los datos que van llegando de los sensores, los datos de los pedidos, incidencias, alertas, ...etc, en resumen, en todos eventos que generan un flujo de datos continuo y que tiene que se procesados de forma rápida, ya sea para actualizar es estado de un pedido o para pedir asistencia en carretera para un vehículo de la flota, o para gestionar una incidencia en una entrega, por ejemplo.
\end{itemize}

\subsection{Equilibrio CAP y SCV}
En esta sección vamos a explicar la aplicación de los teoremas CAP y SCV en nuestro sistema de Big Data, y donde aplicaríamos según que principios.

\begin{itemize}
	\item \textbf{Teorema CAP}: como en todos los sistemas distribuidos, n\textbf{o se puede garantizar 100\%} que no vaya a haber \textbf{particionamiento de la red}, por lo que nos quedaría explicar donde sería mas interesante priorizar la consistencia o la disponibilidad:
	\begin{itemize}
		\item \textbf{Consistencia}: se debería priorizar tanto en las operaciones transaccionales, ya que estas operaciones deben ser consistente y reflejar el estado real de los pedidos, entregas, ..etc, como en los datos históricos del data warehouse, ya que no tenemos especial prisa en procesar estos datos, pero una vez que se procesan deberán ser lo más consistente posibles, para poder generar informes fiables.
		
		\item \textbf{Disponibilidad}: por otro lado, los datos relacionados con la telemetría de los vehículos, así como los datos meteorológicos o de tráfico, ya que nos interesa tener los datos siempre disponibles aunque algunos tengan temporalmente información desactualizada.
	\end{itemize} 
	
	\item \textbf{Teorema SCV}: la aplicación de este teorema es análoga al anterior, que queremos las misma características en los mismos elementos. En este caso, dado nuestro caso, siempre se va a priorizar la \textbf{escalabilidad}, para poder hacer frente al incremento de datos, y su correspondiente almacenamiento y procesamiento. Por lo que nos quedaría elegir lo siguiente:
	
	\item \textbf{Consistencia}: al igual que en el punto anterior, nos interesa tener consistencia en la base de datos relacional, así como en el data warehouse, por lo mismo motivos que en el punto anterior.
	
	\item \textbf{Disponibilidad}: por otro lado en la recepción de los datos telemétricos, sería aceptable que hubiera cierta inconsistencia momentánea, como ya hemos comentado.
\end{itemize}


\subsection{Modelo Predictivo}
Aunque no se ha comentado en los ejercicios anteriores, nuestro sistema incluirá un \textbf{modelo predictivo}, cuya explicación se ha dejado para este punto. 

En nuestro caso, emplearemos un \textbf{modelo de Machine Learning}, ya que no necesitamos un \textbf{LLM} para nuestro propósito. Podemos elegir alguno de los modelos que tenemos en el mercado, como XGBoost o K-Means, por ejemplo. En este caso, se empleará para la \textbf{predicción de retrasos} y \textbf{calculo de rutas}, aunque también lo emplearemos para automatizar un poco las alertas y realizar predicciones sobre ellas. El sistema estaría \textbf{integrado en la capa de análisis}.

Para su \textbf{entrenamiento}, se emplearía un procesamiento por lotes con los datos históricos almacenados en el data warehouse. Es permitiría acelerar el entrenamiento ya que los datos, además de ser consistentes, están preparados para su análisis. 

Una vez entrenado, para su \textbf{funcionamiento diario} se emplearían los datos de telemetría de los vehículos, datos de trafico y meteorológicos, que se usarían para \textbf{predecir} la duración de una ruta, las posibilidades de retraso, detectar anomalías, riesgo de averías de un vehículo, consumo anormal de combustible, ...etc.






\end{document}  